\subsection{贝尔空间上的半连续函数}

定理 2. 设 X 是贝尔空间，\((f_{\alpha})\) 是 X 上一族下半连续实值函数，使得在 X 的每一点 x 处上包络 \(\sup_{\alpha} f_{\alpha}(x)\) 有限。则 X 中每个非空开集都含有一个非空开集，族 \((f_{\alpha})\) 在其上一致有上界。

该定理也可表述成：在其邻域上族 \((f_{\alpha})\) 一致有上界的点所成的集是稠密开集。

设 \(f = \sup_{\alpha} f_{\alpha}\) 是族 \((f_{\alpha})\) 的上包络。函数 \(f\) 是下半连续的（第 IV 章，§ 6，第 2 号，定理 4）且在 \(\mathbf{X}\) 的每一点有限。因此只需就族 \((f_{\alpha})\) 由单独一个函数 \(f\) 组成的情形进行证明。设 \(\mathbf{A}_n\) 为诸点 \(x \in \mathbf{X}\) 中满足 \(f(x) \leqslant n\) 者所成的集；\(\mathbf{A}_n\) 是闭集（第 IV 章，§ 6，第 2 号，命题 1），而诸假设蕴含 \(\mathbf{X}\) 是诸集 \(\mathbf{A}_n\) 的并；故至少有一个 \(\mathbf{A}_n\) 有内点，从而存在一个非空开集，\(f\) 在其上有上界（以某个整数 \(n\) 为界）。若把这个结果应用于 \(\mathbf{X}\) 的任一非空开子集（由第 3 号命题 3 ，该子空间也是贝尔空间），便得定理。

在这定理的应用中，通常 \(f_{\alpha}\) 在 X 上是连续的。

注。若不假设 X 是贝尔空间，定理可能不成立。例如，若对每个有理数 p/q （p、q 为互素的整数，q \textgreater{} 0）令 \(f(p/q) = q\) ，我们就得到有理数直线 Q 上的一个下半连续函数 f ，它在每一点有限（参看第 IV 章，§ 6，第 2 号）；但 Q 中没有使 f 有上界的非空开集。

